\section{Conclusion}\label{sec:conclusion}

In this descriptive, noncausal analysis, most explained movement in the
NMTC's aggregate rural mobilization gap is associated with observed CDE
composition. The order-invariant decomposition associates $-0.185$ of the
raw $-0.262$ gap with CDE identity. The median and two margin specifications
produce small within-CDE estimates that are not statistically distinguishable
from zero. The mean estimate is imprecise and roughly a fifth of the raw gap;
a strictly nested model with state and intermediary effects returns
$-0.060$. Scanning twenty-four subgroups by intermediary size, era, region, project
type, rural orientation and purpose of investment leaves one cell standing
after a correction for scanning: rural commercial real-estate
rehabilitation, at $-0.439$, which survives every transformation, drop,
placebo and resampling test I put to it. That dimension was added after an
earlier draft flagged real estate, so I hold it as exploratory. The cross-CDE distribution of non-metro shares
shows no bunching at the 20\% administrative target. Intermediary composition may itself respond to rural market
conditions, so the mechanism remains unresolved. The estimates motivate
evaluating allocation through the CDE layer alongside project- and
market-level interventions, but they do not establish a causal mechanism or
policy effect.
